\documentclass{amsart}
% For dealing with references we use the comment environment
\usepackage{verbatim}
\newenvironment{reference}{\comment}{\endcomment}
%\newenvironment{reference}{}{}
\newenvironment{slogan}{\comment}{\endcomment}
\newenvironment{history}{\comment}{\endcomment}

% For commutative diagrams we use Xy-pic
\usepackage[all]{xy}

% We use 2cell for 2-commutative diagrams.
\xyoption{2cell}
\UseAllTwocells

% We use multicol for the list of chapters between chapters
\usepackage{multicol}

% This is generally recommended for better output
\usepackage{lmodern}
\usepackage[T1]{fontenc}

% For cross-file-references

% Package for hypertext links:
\usepackage{hyperref}

% For any local file, say "hello.tex" you want to link to please

% Theorem environments.
%
\theoremstyle{plain}
\newtheorem{theorem}[subsection]{Theorem}
\newtheorem{proposition}[subsection]{Proposition}
\newtheorem{lemma}[subsection]{Lemma}

\theoremstyle{definition}
\newtheorem{definition}[subsection]{Definition}
\newtheorem{example}[subsection]{Example}
\newtheorem{exercise}[subsection]{Exercise}
\newtheorem{situation}[subsection]{Situation}

\theoremstyle{remark}
\newtheorem{remark}[subsection]{Remark}
\newtheorem{remarks}[subsection]{Remarks}

\numberwithin{equation}{subsection}

% Macros
%
\def\lim{\mathop{\mathrm{lim}}\nolimits}
\def\colim{\mathop{\mathrm{colim}}\nolimits}
\def\Spec{\mathop{\mathrm{Spec}}}
\def\Hom{\mathop{\mathrm{Hom}}\nolimits}
\def\Ext{\mathop{\mathrm{Ext}}\nolimits}
\def\SheafHom{\mathop{\mathcal{H}\!\mathit{om}}\nolimits}
\def\SheafExt{\mathop{\mathcal{E}\!\mathit{xt}}\nolimits}
\def\Sch{\mathit{Sch}}
\def\Mor{\mathop{\mathrm{Mor}}\nolimits}
\def\Ob{\mathop{\mathrm{Ob}}\nolimits}
\def\Sh{\mathop{\mathit{Sh}}\nolimits}
\def\NL{\mathop{N\!L}\nolimits}
\def\CH{\mathop{\mathrm{CH}}\nolimits}
\def\proetale{{pro\text{-}\acute{e}tale}}
\def\etale{{\acute{e}tale}}
\def\QCoh{\mathit{QCoh}}
\def\Ker{\mathop{\mathrm{Ker}}}
\def\Im{\mathop{\mathrm{Im}}}
\def\Coker{\mathop{\mathrm{Coker}}}
\def\Coim{\mathop{\mathrm{Coim}}}

% Boxtimes
%
\DeclareMathSymbol{\boxtimes}{\mathbin}{AMSa}{"02}

%
% Macros for moduli stacks/spaces
%
\def\QCohstack{\mathcal{QC}\!\mathit{oh}}
\def\Cohstack{\mathcal{C}\!\mathit{oh}}
\def\Spacesstack{\mathcal{S}\!\mathit{paces}}
\def\Quotfunctor{\mathrm{Quot}}
\def\Hilbfunctor{\mathrm{Hilb}}
\def\Curvesstack{\mathcal{C}\!\mathit{urves}}
\def\Polarizedstack{\mathcal{P}\!\mathit{olarized}}
\def\Complexesstack{\mathcal{C}\!\mathit{omplexes}}
% \Pic is the operator that assigns to X its picard group, usage \Pic(X)
% \Picardstack_{X/B} denotes the Picard stack of X over B
% \Picardfunctor_{X/B} denotes the Picard functor of X over B
\def\Pic{\mathop{\mathrm{Pic}}\nolimits}
\def\Picardstack{\mathcal{P}\!\mathit{ic}}
\def\Picardfunctor{\mathrm{Pic}}
\def\Deformationcategory{\mathcal{D}\!\mathit{ef}}

\setlength{\emergencystretch}{3em}
\begin{document}
\title[Stacks Project, Tag 08AG]{582 --- Uniform Castelnuovo\textendash{}Mumford regularity bounds for kernels of fixed-rank trivial-sheaf quotients with fixed Hilbert polynomial}
\maketitle
\noindent Copyright (C) 2005--2025 Johan de Jong. Stacks Project authors. Source: \href{https://stacks.math.columbia.edu/tag/08AG}{Tag 08AG}.

\noindent Editorial difficulty note (not author text). Difficulty H2: Hyperplane induction controls higher cohomology but leaves the first cohomology obstruction unbounded. The proof shows that any nondecreasing step forces all subsequent steps to be injective and hence zero by eventual vanishing, so strict dimension drops and the Hilbert polynomial give a uniform bound..

\noindent Editorial provenance note (not author text). History: excerpt from revision \texttt{a0444\allowbreak 6e57e\allowbreak c1fbc\allowbreak 252a8\allowbreak 71afc\allowbreak ec775\allowbreak 2fb28\allowbreak 07b14}, varieties.tex, lines 6855--6997. Reference presentation was adapted; original-author context and core support blocks were retained or added. The mathematical wording is unchanged. Licensed under GNU FDL 1.2 or later; no invariant sections or cover texts. See ../licenses/COPYING. Context blocks reproduce author definitions and supporting results; they are not additional counted problems.

\begin{definition}
\label{definition-regularity}
Let $k$ be a field. Let $n \geq 0$. Let $\mathcal{F}$ be a coherent
sheaf on $\mathbf{P}^n_k$. We say $\mathcal{F}$ is {\it $m$-regular}
if
$$
H^i(\mathbf{P}^n_k, \mathcal{F}(m - i)) = 0
$$
for $i = 1, \ldots, n$.
\end{definition}

\begin{lemma}
\label{lemma-bound-quotients-free}
Let $k$ be a field. Let $n \geq 0$. Let $r \geq 1$. Let $P \in \mathbf{Q}[t]$.
There exists an integer $m$ depending on $n$, $r$, and $P$
with the following property: if
$$
0 \to \mathcal{K} \to \mathcal{O}^{\oplus r} \to \mathcal{F} \to 0
$$
is a short exact sequence of coherent sheaves on $\mathbf{P}^n_k$
and $\mathcal{F}$ has Hilbert polynomial $P$, then
$\mathcal{K}$ is $m$-regular.
\end{lemma}

\begin{proof}
We prove this by induction on $n$. If $n = 0$, then
$\mathbf{P}^n_k = \Spec(k)$ and any coherent module is $0$-regular
and any surjective map is surjective on global sections.
Assume $n > 0$. Consider an exact sequence as in the lemma.
Let $P' \in \mathbf{Q}[t]$ be the polynomial
$P'(t) = P(t) - P(t - 1)$. Let $m'$ be the integer
which works for $n - 1$, $r$, and $P'$.
By Lemmas \href{https://stacks.math.columbia.edu/tag/08A4}{lemma m regular extend base field (Tag 08A4)} and
\href{https://stacks.math.columbia.edu/tag/08AB}{lemma euler characteristic extend base field (Tag 08AB)}
we may replace $k$ by a field extension, hence we may assume
$k$ is infinite. Apply 
Lemma \href{https://stacks.math.columbia.edu/tag/08A0}{lemma exact sequence induction (Tag 08A0)}
to the coherent sheaf $\mathcal{F}$.
The Hilbert polynomial of $\mathcal{F}' = i^*\mathcal{F}$
is $P'$ (see proof of Lemma \href{https://stacks.math.columbia.edu/tag/08AC}{lemma hilbert polynomial (Tag 08AC)}).
Since $i^*$ is right exact we see that $\mathcal{F}'$ is
a quotient of $\mathcal{O}_H^{\oplus r} = i^*\mathcal{O}^{\oplus r}$.
Thus the induction hypothesis applies to $\mathcal{F}'$ on
$H \cong \mathbf{P}^{n - 1}_k$ (Lemma \href{https://stacks.math.columbia.edu/tag/089Z}{lemma hyperplane (Tag 089Z)}).
Note that the map $\mathcal{K}(-1) \to \mathcal{K}$ is injective
as $\mathcal{K} \subset \mathcal{O}^{\oplus r}$ and has
cokernel $i_*\mathcal{H}$ where $\mathcal{H} = i^*\mathcal{K}$.
By the snake lemma (Homology, Lemma \href{https://stacks.math.columbia.edu/tag/010H}{lemma snake (Tag 010H)})
we obtain a commutative diagram with exact columns and rows
$$
\xymatrix{
& 0 \ar[d] & 0 \ar[d] & 0 \ar[d] \\
0 \ar[r] &
\mathcal{K}(-1) \ar[r] \ar[d] &
\mathcal{O}^{\oplus r}(-1) \ar[r] \ar[d] &
\mathcal{F}(-1) \ar[d] \ar[r] & 0 \\
0 \ar[r] &
\mathcal{K} \ar[r] \ar[d] &
\mathcal{O}^{\oplus r} \ar[r] \ar[d] &
\mathcal{F} \ar[d] \ar[r] & 0\\
0 \ar[r] &
i_*\mathcal{H} \ar[r] \ar[d] &
i_*\mathcal{O}_H^{\oplus r} \ar[r] \ar[d] &
i_*\mathcal{F}' \ar[r] \ar[d] & 0 \\
& 0 & 0 & 0
}
$$
Thus the induction hypothesis applies to the exact sequence
$0 \to \mathcal{H} \to \mathcal{O}_H^{\oplus r} \to \mathcal{F}' \to 0$
on $H \cong \mathbf{P}^{n - 1}_k$ (Lemma \href{https://stacks.math.columbia.edu/tag/089Z}{lemma hyperplane (Tag 089Z)})
and $\mathcal{H}$ is $m'$-regular. Recall that this implies that
$\mathcal{H}$ is $d$-regular for all $d \geq m'$
(Lemma \href{https://stacks.math.columbia.edu/tag/08A6}{lemma m regular up (Tag 08A6)}).

\medskip\noindent
Let $i \geq 2$ and $d \geq m'$. It follows from the long exact
cohomology sequence associated to the left column of the diagram
above and the vanishing of $H^{i - 1}(H, \mathcal{H}(d))$
that the map
$$
H^i(\mathbf{P}^n_k, \mathcal{K}(d - 1))
\longrightarrow
H^i(\mathbf{P}^n_k, \mathcal{K}(d))
$$
is injective. As these groups are zero for $d \gg 0$
(Cohomology of Schemes,
Lemma \href{https://stacks.math.columbia.edu/tag/01YS}{lemma coherent projective (Tag 01YS)})
we conclude $H^i(\mathbf{P}^n_k, \mathcal{K}(d))$ are zero
for all $d \geq m'$ and $i \geq 2$.

\medskip\noindent
We still have to control $H^1$. First we observe that all the maps
$$
H^1(\mathbf{P}^n_k, \mathcal{K}(m' - 1)) \to
H^1(\mathbf{P}^n_k, \mathcal{K}(m')) \to
H^1(\mathbf{P}^n_k, \mathcal{K}(m' + 1)) \to \ldots
$$
are surjective by the vanishing of $H^1(H, \mathcal{H}(d))$ for $d \geq m'$.
Suppose $d > m'$ is such that
$$
H^1(\mathbf{P}^n_k, \mathcal{K}(d - 1))
\longrightarrow
H^1(\mathbf{P}^n_k, \mathcal{K}(d))
$$
is injective. Then
$H^0(\mathbf{P}^n_k, \mathcal{K}(d)) \to H^0(H, \mathcal{H}(d))$
is surjective. Consider the commutative diagram
$$
\xymatrix{
H^0(\mathbf{P}^n_k, \mathcal{K}(d)) \otimes_k
H^0(\mathbf{P}^n_k, \mathcal{O}(1))
\ar[r] \ar[d] &
H^0(\mathbf{P}^n_k, \mathcal{K}(d + 1)) \ar[d] \\
H^0(H, \mathcal{H}(d)) \otimes_k
H^0(H, \mathcal{O}_H(1))
\ar[r] &
H^0(H, \mathcal{H}(d + 1))
}
$$
By Lemma \href{https://stacks.math.columbia.edu/tag/08A7}{lemma m regular multiply (Tag 08A7)}
we see that the bottom horizontal arrow is surjective.
Hence the right vertical arrow is surjective. We conclude that
$$
H^1(\mathbf{P}^n_k, \mathcal{K}(d))
\longrightarrow
H^1(\mathbf{P}^n_k, \mathcal{K}(d + 1))
$$
is injective. By induction we see that
$$
H^1(\mathbf{P}^n_k, \mathcal{K}(d - 1)) \to
H^1(\mathbf{P}^n_k, \mathcal{K}(d)) \to
H^1(\mathbf{P}^n_k, \mathcal{K}(d + 1)) \to \ldots
$$
are all injective and we conclude that
$H^1(\mathbf{P}^n_k, \mathcal{K}(d - 1)) = 0$
because of the eventual vanishing of these groups. Thus the dimensions
of the groups $H^1(\mathbf{P}^n_k, \mathcal{K}(d))$ for $d \geq m'$
are strictly decreasing until they become zero. It follows that the
regularity of $\mathcal{K}$ is bounded by
$m' + \dim_k H^1(\mathbf{P}^n_k, \mathcal{K}(m'))$.
On the other hand, by the vanishing of the higher cohomology groups
we have
$$
\dim_k H^1(\mathbf{P}^n_k, \mathcal{K}(m')) = 
- \chi(\mathbf{P}^n_k, \mathcal{K}(m')) +
\dim_k H^0(\mathbf{P}^n_k, \mathcal{K}(m'))
$$
Note that the $H^0$ has dimension bounded by the dimension of
$H^0(\mathbf{P}^n_k, \mathcal{O}^{\oplus r}(m'))$ which is
at most $r{n + m' \choose n}$ if $m' > 0$ and zero if not.
Finally, the term $\chi(\mathbf{P}^n_k, \mathcal{K}(m'))$ is equal
to $r{n + m' \choose n} - P(m')$. This gives a bound of the
desired type finishing the proof of the lemma.
\end{proof}
\end{document}
